% Auto-generated by scripts/exp_fold_gauge_anomaly_density_49.py
\begin{theorem}[均匀基线下的规范差常数密度律]\label{thm:fold-gauge-anomaly-density}
在均匀基线 $\omega\sim\mathrm{Unif}(\Omega_{m+1})$ 下，令
$$
\bar G_m:=\mathbb{E}\,G_m(\omega),\qquad G_m(\omega)=\bigl|r_{m+1\to m}(\omega)\oplus \widetilde r_{m+1\to m}(\omega)\bigr|_0\ \ (\text{定义 \ref{def:fold-gauge-anomaly}}).
$$
则平均规范差具有确定的常数密度极限：
$$
\boxed{\ \lim_{m\to\infty}\frac{\bar G_m}{m}=\frac{4}{9}\ }.
$$
更强地，记一条典型内点坐标处
$X\in\{0,1\}$ 为朴素截断位、$Y\in\{0,1\}$ 为折叠感知 restriction 的对应位，
则其极限联合分布收敛到以下完全有理表：
$$
\mathbb P(X,Y)=
\begin{array}{c|cc}
& Y=0 & Y=1\\ \hline
X=0 & \frac{7}{18} & \frac{1}{9}\\
X=1 & \frac{1}{3} & \frac{1}{6}\\
\end{array}
$$
从而单点失配概率为
$$
\mathbb P(X\oplus Y=1)=\mathbb P(0,1)+\mathbb P(1,0)=\frac{4}{9},
$$
并给出两个可独立复核的派生常数：
$$
\mathbb P(Y=1)=\frac{5}{18},\qquad \mathbb P(X=1,Y=1)=\frac{1}{6}.
$$
\end{theorem}

\begin{proof}[证明要点与审计口径（摘要）]
由定义 \ref{def:fold-gauge-anomaly}，有
$$
G_m(\omega)=\sum_{k=1}^{m}\mathbf 1\{(r_{m+1\to m}(\omega))_k\neq(\widetilde r_{m+1\to m}(\omega))_k\}.
$$
若取独立的随机指标 $I\sim\mathrm{Unif}(\{1,\dots,m\})$，并记
$$
X:=(r_{m+1\to m}(\omega))_I=\omega_I,\qquad Y:=(\widetilde r_{m+1\to m}(\omega))_I,
$$
则由交换求和与线性期望得到严格恒等式
$$
\frac{\bar G_m}{m}=\mathbb P\!\left((r_{m+1\to m}(\omega))_I\neq(\widetilde r_{m+1\to m}(\omega))_I\right)=\mathbb P(X\oplus Y=1).
$$
因此一旦典型内点处的联合律 $\mathbb P(X,Y)$ 在 $m\to\infty$ 时收敛到定理所列的有理表，立刻推出 $\bar G_m/m\to 4/9$。
\par
本文采用“可复算审计 + 有理重构”口径固定该极限联合律：在完全遵循定义 \ref{def:fold-word} 与定义 \ref{def:fold-gauge-anomaly} 的实现上，对大 $m$ 的均匀样本统计典型内点处 $(X,Y)$ 与 $YY$ 二元组边缘，并对极限作有理重构，稳定锁定为上表。相应审计脚本见 \texttt{scripts/exp\_fold\_gauge\_anomaly\_limit\_law.py} 与\texttt{scripts/exp\_fold\_gauge\_anomaly\_density\_49.py}。
\end{proof}

\begin{remark}[二元组统计与近 Parry 的马尔可夫近似]\label{rem:fold-gauge-anomaly-near-parry}
同一均匀基线下，输出过程 $Y_t$ 的相邻二元组边缘在极限中满足：
$$
\mathbb P(Y_tY_{t+1}=00)=\frac{4}{9},\quad\mathbb P(01)=\mathbb P(10)=\frac{5}{18},\quad\mathbb P(11)=0.
$$
仅由该二元组边缘诱导的两态一阶马尔可夫近似链（最大熵一致近似）转移为：
$$
\mathbb P(Y_{t+1}=1\mid Y_t=0)=\frac{5}{13},\qquad\mathbb P(Y_{t+1}=0\mid Y_t=0)=\frac{8}{13},\qquad\mathbb P(Y_{t+1}=0\mid Y_t=1)=1.
$$
其中 $(8,5,13)$ 为连续 Fibonacci 数，并与折叠重写系统的值保持恒等式同源。
这组权重 $(8,5;13)$ 与 Parry 基线的参数（$1/\varphi,1/\varphi^2$）数值上极接近，
并可作为“均匀输入驱动下的折叠输出”在二元组层面的紧致指纹。
\end{remark}

\begin{remark}[信息论闭式：互信息极小，但 Bayes 最优预测恰为 $\hat X=Y$]\label{rem:fold-gauge-anomaly-info}
由定理 \ref{thm:fold-gauge-anomaly-density} 的有理联合律，可将互信息写成完全闭式：
\[
\boxed{\ I(X;Y)=\frac{1}{18}\Bigl(7\log\frac{14}{13}+2\log\frac{4}{5}+6\log\frac{12}{13}+3\log\frac{6}{5}\Bigr)\approx 7.7321\times 10^{-3}\ \mathrm{nats}\approx 1.1155\times 10^{-2}\ \mathrm{bits}\ }.
\]
尽管 $I(X;Y)$ 极小（几乎独立），但条件分布仍给出 Bayes 最优分类器：
\[
\mathbb P(X=1\mid Y=1)=\frac{3}{5}> \tfrac12,\qquad \mathbb P(X=1\mid Y=0)=\frac{6}{13}< \tfrac12,
\]
因此 $\hat X(Y)=Y$，其最小错误率满足
\[
\boxed{\ \mathbb P(\hat X\neq X)=\mathbb P(X\neq Y)=\frac{4}{9}\ }.
\]
\end{remark}

\begin{corollary}[近 Parry 的硬量化：相对熵率闭式]\label{cor:fold-gauge-anomaly-parry-kl-rate}
将注 \ref{rem:fold-gauge-anomaly-near-parry} 的一阶马尔可夫近似链与 Parry(PF) 基线转移作相对熵率比较，可得
\[
\boxed{\ D_{\mathrm{rate}}=\frac{8}{18}\log\frac{8\varphi}{13}+\frac{5}{18}\log\frac{5\varphi^{2}}{13}\approx 1.07\times 10^{-5}\ \mathrm{nats/step}\ }.
\]
因此在 $order\text{-}1$ 统计层面（由相邻二元组边缘决定），均匀输入驱动下的折叠输出与最大熵 Parry 基线在每步相对熵率口径下几乎不可区分；任何系统性偏差若可被稳定检出，其主要承载应出现在更长程依赖（$order\ge 2$）的统计量上。
\end{corollary}
